\paragrafos{Αντιστρόφως ανάλογα ποσά: περισσότερο το ένα, λιγότερο το άλλο}{Α.6.6, «Αντιστρόφως ανάλογα ποσά»}

\begin{anagke}
Το Σάββατο η Α3 πρέπει να στήσει τα τραπέζια του παζαριού στην αυλή. Ο Στέλιος μετράει: «Πέντε παιδιά, στήσιμο σε $36$ λεπτά». Τότε ο Νίκος λέει «Θα φέρω κι άλλα παιδιά» και ο Στέλιος αναρωτιέται: «Αν διπλασιάσουμε τα παιδιά, θα διπλασιαστεί ο χρόνος;». Όλοι γελούν: με περισσότερα χέρια θα τελειώσουν πιο γρήγορα. Η Μυρτώ παρατηρεί: «Δύο φορές περισσότερα παιδιά, μισός χρόνος».

Η Πολυξένη το γράφει. «Τα πέντε παιδιά κάνουν $36$ λεπτά, δηλαδή $180$ «λεπτά δουλειάς παιδιού». Τα $10$ παιδιά πρέπει να κάνουν κι αυτά $180$ συνολικά, άρα $18$ λεπτά». Ο Βασίλης δοκιμάζει με $9$ παιδιά και βρίσκει $20$, και παρατηρεί ότι το γινόμενο παίδες επί λεπτά βγαίνει πάντα $180$.

Το ίδιο συμβαίνει στο ταξίδι: αν διπλασιάσεις την ταχύτητα, ο χρόνος για την ίδια διαδρομή γίνεται μισός. Και στη μοιρασιά: όσο περισσότερα παιδιά, τόσο λιγότερα μπισκότα το καθένα. Εδώ το ένα ποσό μεγαλώνει και το άλλο μικραίνει με τέτοιον τρόπο, ώστε το γινόμενό τους να μένει ίδιο.
\end{anagke}

\begin{lexi}{αντιστρόφως ανάλογα}{αντιστρόφως ανάλογα}
Τα \lexiword{αντιστρόφως ανάλογα} ποσά είναι αυτά που αλλάζουν «ανάλογα, αλλά ανάποδα» (αντίστροφα): όσες φορές μεγαλώνει το ένα, τόσες μικραίνει το άλλο. Ο \lexiword{αντίστροφος} ενός αριθμού είναι αυτός που, πολλαπλασιασμένος με αυτόν, δίνει $1$.
\end{lexi}

\begin{metafrasi}
Διπλάσια παιδιά, μισός χρόνος. & $x\to2x\ \Rightarrow\ y\to\dfrac{y}{2}$\\
Το γινόμενο των δύο ποσών μένει σταθερό. & $x\cdot y=\alpha$\\
$5$ παιδιά και $36$ λεπτά. & $5\cdot36=180$\\
Τα $10$ παιδιά θα κάνουν. & $y=\dfrac{180}{10}=18$ λεπτά\\
\end{metafrasi}

\tmima{Η θεωρία}

\begin{orismos}{αντιστρόφως ανάλογα ποσά}
Δύο ποσά λέγονται \textbf{αντιστρόφως ανάλογα} όταν μεταβάλλονται με τέτοιον τρόπο ώστε, όταν το ένα πολλαπλασιάζεται επί έναν αριθμό, το άλλο διαιρείται με τον ίδιο αριθμό.
\end{orismos}

Σε αντιστρόφως ανάλογα ποσά $x,\,y$ το \textbf{γινόμενο των αντίστοιχων τιμών} είναι σταθερό:
$$x\cdot y=\alpha\quad(\alpha\neq0)\qquad\text{ή}\qquad y=\dfrac{\alpha}{x}.$$
Στην ειδική περίπτωση $\alpha=1$ τα $x,\,y$ είναι αντίστροφοι αριθμοί ($x\cdot y=1$).

\tmimasmall{Πίνακας και γραφική παράσταση}

Για την ταχύτητα $x$ και τον χρόνο $y$ μιας διαδρομής $120$ km είναι $x\cdot y=120$:
\begin{center}
\renewcommand{\arraystretch}{1.3}
\begin{tabular}{c|cccccc}
ταχύτητα $x$ (km/h) & $30$ & $40$ & $60$ & $80$ & $120$ & $240$\\\hline
χρόνος $y$ (h) & $4$ & $3$ & $2$ & $1{,}5$ & $1$ & $0{,}5$
\end{tabular}
\end{center}
Αν τοποθετήσουμε τα σημεία $(x,\,y)$ σε σύστημα ημιαξόνων, \textbf{δεν} βρίσκονται σε ημιευθεία αλλά πάνω σε μια \textbf{καμπύλη γραμμή}, που πλησιάζει τους άξονες χωρίς να τους ακουμπά.

\begin{center}
\begin{tikzpicture}[x=0.05cm,y=0.095cm]
\draw[gray!35,very thin] (0,0) grid[xstep=30,ystep=1] (250,4.4);
\draw[thick,-{Stealth}] (0,0)--(262,0) node[below] {$x$};
\draw[thick,-{Stealth}] (0,0)--(0,4.9) node[left] {$y$};
\foreach \i in {30,60,120,240}{\node[below,font=\scriptsize] at (\i,0) {\i};}
\foreach \j in {1,2,3,4}{\node[left,font=\scriptsize] at (0,\j) {\j};}
\draw[gblue,thick,domain=28:240,samples=60,smooth] plot (\x,{120/\x});
\foreach \x/\y in {30/4,40/3,60/2,80/1.5,120/1,240/0.5}{\fill[gblue] (\x,\y) circle (2.3pt);}
\end{tikzpicture}
\end{center}

\begin{prosoxi}
\begin{itemize}[leftmargin=1.2em,itemsep=0.3ex,topsep=2pt]
\item Στα αντιστρόφως ανάλογα ποσά, αν το ένα διπλασιάζεται, το άλλο \emph{γίνεται το μισό} (δεν διπλασιάζεται).
\item Δεν ενώνουμε τα σημεία με ημιευθεία και δεν περνά η γραμμή από το $O$: η καμπύλη δεν κόβει τους άξονες (δεν μπορώ να έχω $x=0$ ή $y=0$).
\item Δεν είναι κάθε ζευγάρι ποσών που το ένα μεγαλώνει και το άλλο μικραίνει αντιστρόφως ανάλογο: πρέπει το γινόμενο να μένει σταθερό.
\item Για να δω το είδος της σχέσης, ρωτάω: με διπλάσιο $x$, το $y$ γίνεται διπλάσιο (ανάλογα), μισό (αντιστρόφως ανάλογα) ή τίποτα από τα δύο;
\end{itemize}
\end{prosoxi}

\tmima{Λυμένα παραδείγματα}

\paradeigma{Τα τραπέζια του παζαριού}
\tixero Τα $5$ παιδιά στήνουν τα τραπέζια σε $36$ λεπτά.\par
\tithelo Πόση ώρα θα κάνουν $10$, $9$ και $12$ παιδιά;\par
\possindeo Η δουλειά είναι $5\cdot36=180$ «λεπτά παιδιού». Με $10$: $180:10=18$. Με $9$: $20$. Με $12$: $15$.\par
\apantisi $18$, $20$ και $15$ λεπτά.\par
\elegxos $10\cdot18=9\cdot20=12\cdot15=5\cdot36=180$.

\paradeigma{Η ταχύτητα και ο χρόνος}
\tixero Η διαδρομή είναι $120$ km.\par
\tithelo Ο χρόνος για ταχύτητες $60$, $80$, $48$ km/h.\par
\possindeo $x\cdot y=120$. Για $60$: $2$ ώρες. Για $80$: $1{,}5$. Για $48$: $120:48=2{,}5$.\par
\apantisi $2$, $1{,}5$ και $2{,}5$ ώρες.\par
\elegxos $60\cdot2=80\cdot1{,}5=48\cdot2{,}5=120$.

\paradeigma{Το ορθογώνιο με εμβαδόν $48$}
\tixero Ορθογώνιο με εμβαδόν $48$ m$^2$ και διαστάσεις $x,\,y$.\par
\tithelo Οι ακέραιες διαστάσεις και ποιο έχει τη μικρότερη περίμετρο.\par
\possindeo $x\cdot y=48$: $1\times48$, $2\times24$, $3\times16$, $4\times12$, $6\times8$. Περίμετροι $2(x+y)$: $98$, $52$, $38$, $32$, $28$.\par
\apantisi Η μικρότερη περίμετρος είναι η $28$ m για το $6\times8$ (το πιο «τετράγωνο»).\par
\elegxos Όσο περισσότερο μοιάζει με τετράγωνο, τόσο μικρότερη η περίμετρος για το ίδιο εμβαδόν.

\paradeigma{Είναι αντιστρόφως ανάλογα;}
\tixero Πίνακες: (Α) $x=2,4,6$ και $y=12,6,4$. (Β) $x=2,4,8$ και $y=10,6,3$.\par
\tithelo Ποιος δίνει αντιστρόφως ανάλογα ποσά.\par
\possindeo Γινόμενο: Α: $24$, $24$, $24$. Β: $20$, $24$, $24$.\par
\apantisi Ο Α ναι ($\alpha=24$). Ο Β όχι.\par
\elegxos Στον Α το $x$ διπλασιάζεται από $2$ σε $4$ και το $y$ γίνεται το μισό ($12\to6$).

\paradeigma{Αντίστροφοι αριθμοί}
\tixero Οι αριθμοί $5$, $0{,}5$, $\frac{3}{4}$ και $1{,}25$.\par
\tithelo Οι αντίστροφοί τους.\par
\possindeo Ο αντίστροφος του $x$ είναι ο $y$ με $xy=1$, δηλαδή $y=\frac{1}{x}$. Για το $5$: $\frac{1}{5}$. Για το $0{,}5=\frac{1}{2}$: $2$. Για το $\frac{3}{4}$: $\frac{4}{3}$. Για το $1{,}25=\frac{5}{4}$: $\frac{4}{5}=0{,}8$.\par
\apantisi $\frac{1}{5}$, $2$, $\frac{4}{3}$, $0{,}8$.\par
\elegxos $1{,}25\cdot0{,}8=1$.

\tmima{Ασκήσεις}

\zoni{Το κατάλαβα}{γινόμενα, πίνακες, αντίστροφοι}

\ex Ποιοι πίνακες δίνουν αντιστρόφως ανάλογα ποσά; α) $x$: $2,4,6,12$, $y$: $12,6,4,2$ \quad β) $x$: $1,2,3$, $y$: $5,10,15$ \quad γ) $x$: $3,6,9$, $y$: $8,4,3$ \quad δ) $x$: $5,10,20$, $y$: $8,4,2$.

\ex Τα $x,\,y$ είναι αντιστρόφως ανάλογα με $x\cdot y=36$. Συμπλήρωσε τον πίνακα για $x=1,\,2,\,3,\,4,\,6,\,9$.

\ex Τα $x,\,y$ έχουν $x\cdot y=60$. Βρες το $y$ για $x=5,\ 15,\ 0{,}5$ και το $x$ για $y=30$.

\ex Βρες τους αντίστροφους των $5$, $0{,}5$, $\frac{3}{4}$, $1{,}25$.

\ex Σωστό ή λάθος; α) Σε αντιστρόφως ανάλογα ποσά, όταν το $x$ διπλασιάζεται το $y$ γίνεται το μισό. \quad β) Το γινόμενο $x\cdot y$ είναι σταθερό. \quad γ) Όταν $x=0$, είναι $y=0$. \quad δ) Τα σημεία βρίσκονται σε ημιευθεία από το $O$. \quad ε) Αν $8$ εργάτες κάνουν ένα έργο σε $6$ ημέρες, $12$ εργάτες θέλουν $9$ ημέρες.

\zoni{Το χρησιμοποιώ}{παιδιά, ταχύτητες, δοχεία}

\ex \textbf{Το στήσιμο.} Τα $5$ παιδιά στήνουν το παζάρι σε $36$ λεπτά. Συμπλήρωσε τον πίνακα για $1,\,2,\,3,\,4,\,6,\,9,\,10,\,12,\,15,\,18$ παιδιά.

\ex \textbf{Η ταχύτητα.} Η διαδρομή είναι $120$ km. Βρες τον χρόνο για ταχύτητες $30,\,40,\,48,\,60,\,80,\,120$ km/h.

\ex \textbf{Το ορθογώνιο.} Ένα ορθογώνιο έχει εμβαδόν $48$ m$^2$. Βρες τις ακέραιες διαστάσεις του και πες ποιο έχει τη μικρότερη περίμετρο.

\ex \textbf{Οι βρύσες.} Τέσσερις ίδιες βρύσες γεμίζουν μια δεξαμενή σε $6$ ώρες. Σε πόσες ώρες θα τη γεμίσουν $3$ βρύσες και σε πόσες $8$;

\ex \textbf{Τα μπισκότα.} Μοιράζουμε $72$ μπισκότα εξίσου σε $x$ παιδιά. Βρες πόσα παίρνει το καθένα για $x=8,\ 12,\ 18$.

\ex \textbf{Το διάβασμα.} Το βιβλίο έχει $60$ σελίδες. Πόσες ώρες θέλει η Ζωή να το διαβάσει με $20$ σελίδες ανά ώρα και πόσες με $30$ σελίδες ανά ώρα; Πόσες σελίδες την ώρα πρέπει να διαβάζει για να το τελειώσει σε $1{,}5$ ώρες;

\ex \textbf{Το ταμείο.} Τα $40$ € του ταμείου μοιράζονται εξίσου σε $8$ παιδιά και σε $10$ παιδιά. Πόσο παίρνει το καθένα στην κάθε περίπτωση;

\zoni{Το σκέφτομαι}{ανάλογα ή αντιστρόφως;}

\ex Ο Στέλιος λέει «διπλάσια παιδιά, διπλάσιος χρόνος». Πού κάνει λάθος;

\ex Τοποθέτησε στο ίδιο σύστημα αξόνων τα σημεία $(1,\,48)$, $(2,\,24)$, $(3,\,16)$, $(4,\,12)$, $(6,\,8)$, $(8,\,6)$, $(12,\,4)$ και ένωσέ τα. Είναι ημιευθεία; Φτάνει η γραμμή στους άξονες;

\ex Αν η ταχύτητα διπλασιαστεί, τι γίνεται ο χρόνος για την ίδια διαδρομή; Και αν τριπλασιαστεί;

\ex Πες για κάθε ζευγάρι αν τα ποσά είναι ανάλογα, αντιστρόφως ανάλογα ή τίποτα από τα δύο: α) κιλά και τιμή \quad β) ταχύτητα και χρόνος \quad γ) ηλικία και ύψος \quad δ) αριθμός ατόμων και μερίδιο από σταθερό ποσό \quad ε) πλευρά και εμβαδόν τετραγώνου.

\ex Αν τα $x,\,y$ έχουν $xy=36$, δείξε ότι τα $y$ και $\frac{1}{x}$ είναι ανάλογα και βρες τον συντελεστή.

\zoni{Για το σπίτι}{πόσο γρήγορα φτάνω;}

\ex Μέτρησε την απόσταση από το σπίτι ως το σχολείο και υπολόγισε τον χρόνο για ταχύτητες περπατήματος, ποδηλάτου και αυτοκινήτου.

\ex Βρες δύο παραδείγματα αντιστρόφως αναλόγων ποσών από την καθημερινή ζωή.
